\section{กฎคำร้องและกระบวนการอนุมัติ}
\subsection{เงื่อนไขของคำร้อง}
\begin{itemize}
  \item ประเภทการลาคือ ลาป่วย หรือ ลากิจ
  \item ต้องกรอกเหตุผล วันที่เริ่ม วันที่สิ้นสุด และเลือกรายวิชาอย่างน้อยหนึ่งวิชา
  \item วันสิ้นสุดต้องไม่ก่อนวันเริ่ม และจำนวนวันทำการต้องมากกว่าศูนย์
  \item เบอร์โทรศัพท์ที่กรอกใหม่ต้องมี 10 หลัก
  \item หากไม่ใช้ที่อยู่เดิม ต้องกรอกรายละเอียดและเลือกรหัสไปรษณีย์
  \item ลาป่วยบังคับแนบหลักฐาน รองรับ JPG, JPEG, PNG, WEBP และ PDF ไม่เกิน 2 MB
\end{itemize}

\subsection{เส้นทางอนุมัติ}
เมื่อสร้างคำร้อง ระบบอ่านข้อมูลนิสิต อาจารย์ที่ปรึกษา และโครงสร้างภาควิชา จากนั้นสร้างเส้นทางอนุมัติตามลำดับดังนี้
\begin{enumerate}
  \item \role{ADVISOR}: อาจารย์ที่ปรึกษา
  \item \role{DEPARTMENT\_HEAD}: หัวหน้าภาควิชา
  \item \role{ASSOCIATE\_DEAN}: รองคณบดี เฉพาะเมื่อจำนวนวันลามากกว่า 15 วัน
\end{enumerate}

หากไม่พบผู้อนุมัติแม้แต่หนึ่งราย ระบบจะไม่สร้างคำร้องและแสดงข้อความต่อไปนี้
\begin{center}
  \code{No approval route found for this student}
\end{center}

\subsection{สถานะ}
\begin{table}[H]\centering
  \begin{tabularx}{\textwidth}{|>{\ttfamily\fontsize{13pt}{16pt}\selectfont\raggedright\arraybackslash}p{3.2cm}|X|}\hline
    \tableheader{สถานะ} & \tableheader{ความหมาย} \\\hline
    PENDING & คำร้องกำลังรอการพิจารณา \\\hline
WAITING & ขั้นยังทำรายการไม่ได้เพราะรอขั้นก่อน \\\hline
CURRENT & ขั้นปัจจุบัน ผู้มีสิทธิ์พิจารณาได้ \\\hline
APPROVED & ขั้นหรือคำร้องผ่านแล้ว \\\hline
REJECTED & ไม่อนุมัติและ workflow สิ้นสุด \\\hline
  \end{tabularx}
  \caption{สถานะคำร้องและขั้นอนุมัติ}
\end{table}
การเปลี่ยนสถานะเป็นไปตามกฎต่อไปนี้
\begin{itemize}
  \item ขั้นแรกเริ่มต้นด้วยสถานะ \code{CURRENT}
  \item เมื่ออนุมัติ ระบบจะเลื่อนขั้นถัดไปเป็น \code{CURRENT}
  \item หากไม่มีขั้นถัดไป คำร้องจะเป็น \code{APPROVED}
  \item หากขั้นใดไม่อนุมัติ คำร้องจะเป็น \code{REJECTED} ทันที
\end{itemize}

\subsection{การอนุมัติอัตโนมัติ}
ระบบมีกฎอนุมัติอัตโนมัติแยกตาม role และงานตามเวลา ผลมีสถานะเทียบเท่าการอนุมัติโดยผู้ใช้ จึงต้องกำหนดตามนโยบาย ทดสอบแยกจาก production และติดตาม log หลังเปลี่ยนค่า
